\section{Root-finding and optimization}
\thispagestyle{plain}

\subsection{Overview on numerical methods for minimizing the loss (in the MLE setting)}
The basic methods are to either numerically minimize the loss function directly, e.g.
\begin{itemize}
    \item gradient descent
\end{itemize}
or to apply root-finding methods to the gradient of the loss function, in the MLE
setting the gradient of the log-likelihood, \textbf{the score}
\begin{equation}
    \vec{s}(\vec{\theta}) = \nabla_{\vec{\theta}} \log \mathcal{L}(\vec{\theta})
\end{equation}
with methods like
\begin{itemize}
    \item Newton-Raphson
    \item Fisher scoring\footnote{Like Newton-Raphson but with the Fisher information matrix in place of the Hessian.
    \begin{equation}
        \mat{I}(\vec{\theta}) = E(\vec{s} \vec{s}^T) =-E \left( \frac{\partial^2 \log \mathcal{L}}{\partial \vec{\theta} \partial \vec{\theta}^T} \right)
    \end{equation}
    (the negative expectation of the Hessian of the log-likelihood)
    (which is not a function of any particular observation as the random variable $X$ was 
    averaged out).}
\end{itemize}

\subsection{Gradient Descent}
Here we just walk down the gradient of the loss function, starting at some initial guess $\vec{\theta}_0$, so
\begin{equation}
    \begin{gathered}
        \vec{\hat{\theta}}^{(n+1)}=\vec{\hat{\theta}}^{(n)}-\left.\gamma \vec{\nabla}_{\vec{\theta}} L\right|_{\vec{\theta}=\vec{\widehat{\theta}}^{(n)}}, \quad \text{learning rate } \gamma \\
        \text{until } \left|L\left(\vec{\hat{\theta}}^{(n+1)}\right)-L\left(\vec{\hat{\theta}}^{(n)}\right)\right|<\epsilon \text{ or \# steps } >N_{\text {max }}
    \end{gathered}
\end{equation}
where $\epsilon$ is a small number and $N_{\text{max}}$ is the maximum number of steps, set by
the user. Alternatively, one can stop when $||\vec{\theta}^{(n+1)} - \vec{\theta}^{(n)}|| < \epsilon$.

\redbox{\textbf{Problems of standard gradient descent:}
\begin{itemize}
    \item finds local minima
    \item how to choose the learning rate
    \item gets stuck on plateaus
\end{itemize}
}

\subsubsection{Improvements to standard gradient descent}
\begin{itemize}
    \item to combat local minima, we can start from multiple initial conditions and take the best final result
    \item we prefer an approximate global optimum over an exact local one $\rightarrow$ use stochasticity, 
    (e.g. injected noise $\rightarrow$ jump out of local but not global minima); simulated annealing (probabilistically 
    move to some state or stay in the current one, even locally suboptimal steps are allowed)
    \item \textbf{stochastic gradient descent}: consider we want to minimize an additive function
    \begin{equation}
        Q(\vec{\theta}) = \frac{1}{N} \sum_{i=1}^N Q_i(\vec{\theta}), \quad \text{e.g.} L(\vec{\theta}) = -\mathcal{LL}(\vec{\theta}) = \sum_{i=1}^N -\log p(\vec{x}_i | \vec{\theta})
    \end{equation}
    We do the gradient descent steps by
    \begin{equation}
        \vec{\hat{\theta}}^{(n+1)}=\vec{\hat{\theta}}^{(n)}-\frac{\gamma}{k} \vec{\nabla}_{\vec{\theta}} \sum_{\substack{\text { random subset } \\ \text { of size } k}} Q_k(\vec{\theta})
    \end{equation}
    so we only consider the loss following from a subset of data points - which can avoid local minima (noise
    in the data itself is used)
    \item adaptive learning rage $\gamma$ (reduce for steep gradients) (e.g. AdaGrad, RMSProp, Adam)
\end{itemize}

\subsection{Newton-Raphson}

\subsubsection{Standard Newton iteration for root finding}
Consider we want to find the root of a function $g(\xi)$. We linearly approximate
$g$ around the current guess $\xi_k$ and solve for the root.
\begin{equation}
    \begin{gathered}
    g\left(\xi_{k+1}\right) \approx g\left(\xi_k\right)+\left.\left(\xi_{k+1}-\xi_k\right) \partial_{\xi} g\right|_{\xi=\xi_k}=0 \\
    \rightarrow \xi_{k+1}=\xi_k-\frac{g\left(\xi_k\right)}{\left.\partial_{\xi} g\right|_{\xi=\xi_k}}
    \end{gathered}
\end{equation}
so in the multivariate case, $g: \mathbb{R}^n \rightarrow \mathbb{R}^m$, we have
\begin{equation}
    \vec{\xi}_{k+1}=\vec{\xi}_k-\mat{J}_{\vec{g}}^{-1}\left(\vec{\xi}_k\right) \vec{g}\left(\vec{\xi}_k\right), \quad \mat{J} \vec{g}=\left(\begin{array}{ccc}
    - & \vec{\nabla}_{\vec{\xi}} g_1 & - \\
    \vdots & \vdots & \vdots \\
    - & \vec{\nabla}_{\vec{\xi}} g_n & -
    \end{array}\right)
\end{equation}
Which is illustrated in Fig. \ref{fig:newton_raphson}.

\begin{figure}[!htb]
    \centering
    \includesvg[width=0.7\textwidth]{figures/newton_raphson.svg}
    \caption{Illustration of the Newton-Raphson method.}
    \label{fig:newton_raphson}
\end{figure}

\subsubsection{Newton's method in optimization}
We now want to find critical values $\xi_{\text{crit}}$ of a function $g$, where
\begin{equation}
    \left. \partial_{\xi} g \right|_{\xi_{\text{crit}}} = 0
\end{equation}
\redbox{These can be minima, maxima or saddle points.}
Using Newton iteration, we get
\begin{equation}
    \xi_{k+1}=\xi_k-\frac{\left.\partial_{\xi} g\right|_{\xi=\xi_k}}{\left.\partial_{\xi}^2 g\right|_{\xi=\xi_k}}
\end{equation}
\greenbox{\textbf{Intuition of Newton-optimization:} At
$\xi_k$, a parabola is fitted to the function $g$ and we move the 
the extremum (in higher dimension this can be a saddle point) of it.}

In higher dimensions $g: \mathbb{R}^n \rightarrow \mathbb{R}$, we get
\begin{equation}
    \vec{\xi}_{k+1} = \vec{\xi}_k - \mat{H}_{\vec{g}}^{-1}(\vec{\xi}_k) \vec{\nabla}_{\vec{\xi}} g(\vec{\xi}_k)
\end{equation}
with the numerically costly Hessian
\begin{equation}
    \mat{H}_g = \mat{J}_{\vec{\nabla}_{\vec{\xi}} g}
\end{equation}
best analytically found or calculated based on Automatic 
Differentiation (AD).
\greenbox{In our loss-optimization we simply have $g(\vec{\xi}) = L(\vec{\theta})$.}

\subsubsubsection{Advantages of Newton optimization}
\begin{itemize}
    \item In the convex / concave setting the usage of curvature information
    by Newton iteration leads to faster convergation than gradient descent, as
    illustrated in Fig. \ref{fig:newton_vs_gd}.
    \item step-size is automatically adaptive
    \begin{itemize}
        \item \textit{small Hessian} $\rightarrow$ automatically larger steps
        \item \textit{large Hessian} $\rightarrow$ automatically smaller steps
    \end{itemize}
\end{itemize}

\begin{figure}[!htb]
    \centering
    \includesvg[width=0.4\textwidth]{figures/newton_vs_gd.svg}
    \caption{Newton optimization (\textcolor{green1}{green}), gradient descent (\textcolor{red1}{red}).}
    \label{fig:newton_vs_gd}
\end{figure}

\subsubsubsection{Disadvantages of Newton optimization}
\begin{itemize}
    \item strictly only applies to convex / concave functions $g$
    \item at the inflection point of a non-convex function, the Hessian is indefinite
    so in the 1D case $\partial_{\xi}^2 g = 0$ and Newton optimization breaks down
    \item it searches critical points, not minima, if $\partial_{\xi}^2 g < 0$ and $\partial_{\xi} g > 0$,
    we go uphill, see Fig. \ref{fig:newton_crit}.
\end{itemize}

\begin{figure}[!htb]
    \centering
    \includesvg[width=0.6\textwidth]{figures/newton_crit.svg}
    \caption{Newton optimization breaks down at the inflection point and goes uphill right of it.}
    \label{fig:newton_crit}
\end{figure}

\pagebreak